%%%%%%%%%%%%%%%%%%%%%%%%%%%%%%%%%%%%%%%%%%%%%%%%%%%%%%%
% % Lines starting with % are comments, which are ignored.
% % This is a handy way of indicating the date and version of
% % your document, to wit:
% %
% % LaTeX sample file
% % Modified March, 2002
% %
%%%%%%%%%%%%%%%%%%%%%%%%%%%%%%%%%%%%%%%%%%%%%%%%%%%%%%%
\DocumentMetadata{
  % PDF standards: ONE key whose value is a comma-separated list --
  % UA-2 (accessibility) plus A-4f (archival; the archival companion of
  % UA-2). Do NOT use separate pdfstandard=... keys: with ordinary
  % key-value parsing only the last one takes effect.
  pdfstandard={UA-2,A-4f},
  lang=en-US,
  tagging=on,
  tagging-setup={
    math/setup=mathml-SE
  }
}

\documentclass{ltx-talk}
% lecture-macros.tex -- shared package setup for lecture-template.tex
% (ltx-talk slides / handouts).
%
% NOTE ON PATHS: every template in this repository is stored one level up
% from where it is compiled (see README.md). All relative \input paths
% below are resolved against the folder the document is compiled in,
% which sits one directory above this file.

% --- Tables, bookmarks, strikethrough ------------------------------------
% booktabs: professional table rules (\toprule, \midrule, \bottomrule)
% bookmark: PDF outline entries (PDF bookmarks)
% ulem[normalem]: \sout strikethrough that survives math mode
\usepackage{booktabs}
\usepackage{bookmark}
\usepackage[normalem]{ulem}
\usepackage{multicol}

% --- Math ------------------------------------------------------------
% amsmath: equation environments and numbering
% mathrsfs: \mathfrak  |  stmaryrd: extra symbols
\usepackage{amsmath}
\usepackage{mathrsfs}
\usepackage{stmaryrd}

% --- Layout -------------------------------------------------------------
% float: [H] float placement
% fontspec: Unicode/OpenType fonts (compile with XeLaTeX or LuaLaTeX)
\usepackage{float}
\usepackage{fontspec}

% --- Colors ---------------------------------------------------------------
% dvipsnames: extended named color palette (the Accessible* colors below
% are built on it);
% table: \rowcolors / \cellcolor for shading tables.
% NOTE: xcolor may only be loaded ONCE per document -- a second
% \usepackage{xcolor} with different options is a fatal "Option clash"
% error, so ALL of its options are listed here.
\usepackage[dvipsnames,table]{xcolor}

% --- Multimedia -----------------------------------------------------------
% animate: embeds audio/video (and legacy Flash) media in a frame.
\usepackage{animate}

% --- Bibliography ----------------------------------------------------------
% biblatex (biber backend): citations for lecture notes.
\usepackage[
  backend=biber,
  style=alphabetic,
  url=true,
  doi=true,
  isbn=true,
  backref=true
]{biblatex}

% --- Colored boxes ----------------------------------------------------------
% talkthemetcolorbox: the colored-box theme for ltx-talk slides
% (shadow/rounded/shaded blocks, showtitle, titlepage styles).
\usepackage[
  shadow,
  rounded,
  shaded,
  blocks,
  showtitle,
  titlepage
]{talkthemetcolorbox}

% --- "theindex" helper environment ------------------------------------------
% Provides a lightweight theindex environment (an itemize) and the
% \@idxitem/\subitem/\subsubitem/\indexspace macros that LaTeX's standard
% theindex machinery expects, so hand-written index pages indent correctly.
\makeatletter
\newenvironment{theindex}
  {\begin{itemize}}
  {\end{itemize}}
\providecommand\@idxitem{\par\hangindent 40\p@}
\providecommand\subitem{\@idxitem \hspace*{20\p@}}
\providecommand\subsubitem{\@idxitem \hspace*{30\p@}}
\providecommand\indexspace{\par \vskip 10\p@ \@plus5\p@ \@minus3\p@\relax}
\makeatother

% --- Index -----------------------------------------------------------------
% imakeidx: builds the printed index.
% The redefinition below appends "|hyperpage" to every index entry so that,
% in the PDF, each index entry is a link to its page.
\usepackage{imakeidx}
\makeatletter
\let\imki@orig@wrindexentrysplit\imki@wrindexentrysplit
\renewcommand\imki@wrindexentrysplit[3]{%
  \imki@orig@wrindexentrysplit{#1}{#2|hyperpage}{#3}%
}

\makeatletter
% --- Section titles: capture + one bookmark per logical section -------------
% \section is redefined so that:
%   1. the section title is captured into the plain-TeX register
%      \mytalksectitletoks (reprinted verbatim on the auto divider frame,
%      see the cmd/section/after hook below -- plain-TeX storage is
%      catcode-safe, unlike expanding the title in a hook), and
%   2. exactly ONE depth-1 PDF bookmark is created per \section call
%      (hyperref's own bookmark is suppressed via \Hy@writebookmark).
\DeclareCommandCopy{\OriginalSectionForBookmarks}{\section}
\newtoks\mytalksectitletoks

\RenewDocumentCommand{\section}{s o m}{%
  \begingroup
    \def\Hy@writebookmark##1##2##3##4##5{}%
    \global\mytalksectitletoks={#3}%
    \IfBooleanTF{#1}
      {\phantomsection\IfNoValueTF{#2}{\OriginalSectionForBookmarks*{#3}}{\OriginalSectionForBookmarks*[#2]{#3}}}%
      {\IfNoValueTF{#2}{\OriginalSectionForBookmarks{#3}}{\OriginalSectionForBookmarks[#2]{#3}}}%
  \endgroup
  \IfBooleanTF{#1}
    {\IfNoValueTF{#2}{\pdfbookmark[1]{#3}{section-star.\thepage.\number\inputlineno}}{\pdfbookmark[1]{#2}{section-star.\thepage.\number\inputlineno}}}%
    {\IfNoValueTF{#2}{\pdfbookmark[1]{Section \thesection: #3}{section.\thesection}}{\pdfbookmark[1]{Section \thesection: #2}{section.\thesection}}}%
}
\makeatother

% --- One PDF bookmark per logical frame (NOT per overlay) ---------------------
% ltx-talk writes one PDF page per overlay, so a naive frametitle hook would
% register a new bookmark for every overlay of the same frame. The
% \g_slideutils_frame_bookmark_int register remembers the last frame that
% got a bookmark, so subsequent overlays of the same frame are skipped.
% This is also what keeps BeamerPresenter's cumulative annotations working:
% one logical frame must stay one logical unit in the PDF outline.
\ExplSyntaxOn
\int_new:N \g_slideutils_frame_bookmark_int

\AddToHook{cmd/frametitle/after}{%
  \tl_if_empty:NF \g__talk_frame_title_tl
    {
      \ifnum\value{frame}=\g_slideutils_frame_bookmark_int\relax
      \else
        \int_gset:Nn \g_slideutils_frame_bookmark_int { \value{frame} }%
        \pdfbookmark[2]{\tl_use:N \g__talk_frame_title_tl}
          {frame.\int_use:N \g_slideutils_frame_bookmark_int}
      \fi
    }
}
\ExplSyntaxOff
\makeatother

% --- Printed index + theme for the colored boxes -------------------------------
% \makeindex[...]: two-column index titled "Index", listed in the TOC.
% The two EditInstance calls style the standard header/footer instances
% of talkthemetcolorbox (cyan background; footer order title/author/
% date/framenumber).
\makeindex[columns=2, title=Index, intoc]\EditInstance{header}{std}{
  background-color = cyan,
  color            = black,
  height           = 0.75cm
}
\EditInstance{footer}{std}{
  background-color = cyan,
  element-order = {title, author, date, framenumber},
}

% --- Auto section divider frame ---------------------------------------------------
% After each \section, insert a centered, large, structure-colored frame
% that prints the section title (from \mytalksectitletoks).
\AddToHook{cmd/section/after}{%
  \begin{frame}
    \centering
    \color{structure}
    \LARGE
    \the\mytalksectitletoks
  \end{frame}%
}

% --- Tagging sockets (tagged PDF structure) ---------------------------------------
% ltx-talk exposes tagging "sockets" where the PDF structure can attach
% section titles; register one and leave it unassigned (none) so the
% class's own tagging keeps control.
\ExplSyntaxOn
\AtBeginDocument{
  \NewTaggingSocketPlug { talk / sec / title } { none } { }
  \AssignTaggingSocketPlug { talk / sec / title } { none }
}
\ExplSyntaxOff

% --- Accessibility roles ----------------------------------------------------------------
% Map \frametitle to the H1 heading role in the tagged PDF, so screen
% readers treat slide titles as headings.
\tagpdfsetup{role / new-tag = frametitle / H1}

\makeatletter
% --- Page label = frame number -------------------------------------------------------
% Redefine what feeds the PDF page label to track the frame counter,
% not the physical page counter. Replace \theframenumber below with
% whatever internal macro/property ltx-talk actually uses to produce
% the "framenumber" footer element — confirm via ltx-talk-frame.dtx.
\renewcommand\thepage{\arabic{frame}}
\makeatother

\addbibresource{refs.bib}

% --- "Problem Time!" slides ------------------------------------------------------------------
% \ProblemSlide{problem}{solution} emits two frames: the problem, then
% its solution after a \pause, each numbered by the problem counter.
\newcounter{problemcounter}
\setcounter{problemcounter}{0}
\newcommand{\ProblemSlide}[2]{
  \stepcounter{problemcounter}
  \begin{frame}
    \frametitle{Problem Time!}
    Problem \theproblemcounter: #1
  \end{frame}

  \begin{frame}
    \frametitle{Problem Time!}
    Solution to Problem \theproblemcounter: \pause

    #2
  \end{frame}
}
% \highlight[<colour>]{<stuff>}
% from https://tex.stackexchange.com/questions/33401/a-version-of-colorbox-that-works-inside-math-environments
\newcommand{\highlight}[2][yellow]{\mathchoice%
  {\colorbox{#1}{$\displaystyle#2$}}%
  {\colorbox{#1}{$\textstyle#2$}}%
  {\colorbox{#1}{$\scriptstyle#2$}}%
  {\colorbox{#1}{$\scriptscriptstyle#2$}}}%

% --- Fonts ------------------------------------------------------------------------------
% Lato (text) + LeteSansMath (math): open, accessible typefaces.
\setmainfont{Lato}[
  Ligatures=TeX,
  Numbers={Lining,Monospaced}
]
\setsansfont{Lato}[
  Ligatures=TeX,
  Numbers={Lining,Monospaced}
]
\setmathfont{LeteSansMath.otf}\usepackage{emoji}

% --- Accessible colors ---------------------------------------------------------------------
% WCAG-contrast-safe variants of red/green for text on white:
% red!80!black and green!40!black both clear the 4.5:1 requirement.
% AltYellow: soft background for highlight boxes.
\colorlet{AccessibleRed}{red!80!black}
\colorlet{AccessibleGreen}{green!40!black}
\definecolor{AltYellow}{RGB}{255,245,170}

% --- Links ---------------------------------------------------------------------------------
% hyperref + link colors: every link uses AccessibleGreen (7.2:1 on
% white), so internal, citation, and URL links all pass WCAG 1.4.3.
\usepackage{hyperref}
\hypersetup{
  colorlinks=true,
  linkcolor=AccessibleGreen,
  urlcolor=AccessibleGreen,
  citecolor=AccessibleGreen,
}

% --- Shared course macros ---------------------------------------------------------------------
% Root macros.tex: math shorthands used across all course documents
% (\RR, \NN, \ZZ, \QQ, \grad, \curl, \tr, \der, \pder, ...).
\newcommand{\RR}{\mathbb{R}}
\newcommand{\NN}{\mathbb{N}}
\newcommand{\ZZ}{\mathbb{Z}}
\newcommand{\QQ}{\mathbb{Q}}
\newcommand{\vv}{\vec{v}}
\newcommand{\ww}{\vec{w}}
\newcommand{\xx}{\vec{x}}
\newcommand{\yy}{\vec{y}}
\newcommand{\zz}{\vec{z}}
\newcommand{\uu}{\vec{u}}
\newcommand{\basis}[1]{\vec{e_{#1}}}
\newcommand{\id}{\mathrm{id}}
\newcommand{\CC}{\mathbb{C}}
\renewcommand{\div}{\operatorname{div}}
\newcommand{\grad}{\operatorname{grad}}
\newcommand{\curl}{\operatorname{curl}}
\newcommand{\tr}{\operatorname{tr}}
\newcommand{\der}[2]{\dfrac{d #1}{d #2}}
\newcommand{\pder}[3][]{\dfrac{\partial^{#1} #2}{\partial #3}}


% % Title and author(s)
%%%%%%%%%%%%%%%%%%%%%%%%%%%%%%%%%%%%%%%%%%%%%%%%%%%%%%%
\title{Lecture 13: Related Rates I}
\author{Garrett Credi}
\date{10/22/2026}
%%%%%%%%%%%%%%%%%%%%%%%%%%%%%%%%%%%%%%%%%%%%%%%%%%%%%%%
% %
% % The next command allows your in import encapsulated
% % postscript files, .epsf or .eps files, which
% % contain vector graphic image data.
% %
%%%%%%%%%%%%%%%%%%%%%%%%%%%%%%%%%%%%%%%%%%%%%%%%%%%%%%%
% %   The first thanks indicates your affiliation
% %
% %  Just the name here.
% %
% % Your mailing address goes at the end.
% %
% % \thanks is also how you indicate grant support
% %
%%%%%%%%%%%%%%%%%%%%%%%%%%%%%%%%%%%%%%%%%%%%%%%%%%%%%%%

%%%%%%%%%%%%%%%%%%%%%%%%%%%%%%%%%%%%%%%%%%%%%%%%%%%%%%%
% $ The next command ensures the margins are 1 inch on
% % all sides.
% %
%%%%%%%%%%%%%%%%%%%%%%%%%%%%%%%%%%%%%%%%%%%%%%%%%%%%%%%
\begin{document}
\begin{frame}
  \maketitle
\end{frame}

\begin{frame}
  \frametitle{Exam 2 Soon!}

  Exam 2 is coming up! It is two Tuesday's from now, 11/3.

  It's covering linearizations + our derivative rules + related rates + L'Hôpital's Rule
  which we'll cover next week.
\end{frame}

\begin{frame}
  \frametitle{Quiz Time}

  Let's get into neat rows, spread out, and slay this quiz!
\end{frame}

\begin{frame}
  \frametitle{Implicit Differentiation Recap}

  Last lecture we covered \textit{implicit differentiation},
  the process of relating variables using some implicit relation \pause
  and differentiating that relation to get a relation between their derivatives. \pause

  E.g.: slope of line tangent to $y^2 = x^3 - 10x + 1$ at $(0,1)$. \pause

  Slope of tangent line $\implies$ interested in $\dfrac{dy}{dx}$ $\implies$ treat $y$ as a function of $x$. \pause

  \begin{gather}
    y^2 = x^3 - 10x + 1 \\ \pause
    \dfrac{d}{dx} y^2 = \dfrac{d}{dx} (x^3 - 10x + 1) \\ \pause
    \dfrac{dy}{dx} \dfrac{d}{dy} y^2 = 3x^2 - 10 \\ \pause
    2y y' = 3x^2 - 10 \\ \pause
    2 \cdot 1 \cdot y' = 3 \cdot 0^2 - 10 \\ \pause
    2 y' = -10 \pause \quad y' = -5 
  \end{gather}
\end{frame}

\begin{frame}
  \frametitle{Towards Related Rates}

  Suppose instead that some particle was moving over time in the $2D$ plane and whose motion was constrained to 
  that curve cut out by $y^2 = x^3 - 10x + 1$. \pause
  \begin{itemize}
    \item Moving over time $\implies$ it's position can be described by functions of a time variable $t$ \pause
    \item In the $2D$ plane $\implies$ it has an $x$-position $x(t)$ and a $y$-position $y(t)$ \pause
    \item Motion constrained $\implies$ $y(t)^2 = x(t)^3 - 10x(t) + 1$. \pause
  \end{itemize}

  In order to stay on that curve, what restrictions does that place on the particle's: \pause
  \begin{enumerate}
    \item Position? \pause
    \item Velocity? \pause
    \item Acceleration? 
  \end{enumerate}
\end{frame}

\begin{frame}
  \frametitle{Relating Some Rates I}
  Let's see what our constraint on position forces upon our velocity! \pause

  \begin{gather}
    y(t)^2 = x(t)^3 - 10x(t) + 1 \\ \pause
    \text{Differentiate w.r.t. } t \text{ to get velocity} \\ \pause
    \dfrac{d}{dt} y(t)^2 = \dfrac{d}{dt}(x(t)^3 - 10x(t) + 1) \\ \pause
    \dfrac{dy}{dt} \dfrac{d}{dy} y^2 = \dfrac{dx}{dt} \dfrac{d}{dx}(x^3 - 10x + 1) \\ \pause
    2y\dfrac{dy}{dt} = (3x^2 - 10)\dfrac{dx}{dt}
  \end{gather}\pause
\end{frame}
\begin{frame}
  \frametitle{Relating Some Rates II}
  So if the particle is at $(0,1)$ at some time $t = t_0$ and has a horizontal velocity of $3 m/s$ \pause
  which means $(x(t_0),y(t_0)) = (0,1)$ and $\dfrac{dx}{dt}\mid_{t = t_0} = 3$ what does that force the vertical velocity to be? \pause

  \begin{gather}
    2y \dfrac{dy}{dt} = (3x^2-10)\dfrac{dx}{dt} \\ \pause
    2 \cdot 1 \dfrac{dy}{dt} = (3 \cdot 0^2 - 10) \cdot 3 \\ \pause
    2 \dfrac{dy}{dt} = -30 \quad \pause \dfrac{dy}{dt} = -15 m/s
  \end{gather} \pause

  (Notice from before, we found the slope of the line tangent to the curve at $(1,0)$ was $-5$ and $-15 = -5 \cdot 3$,
  so we're exactly enforcing that the velocity is parallel to the tangent line).
\end{frame}

\begin{frame}
  \frametitle{Related Rates I}

  This process is common enough that we call problems using it \textit{related rates} problems: \pause

  \begin{block}{Related Rates}\index{related rates}
    Given a real-world/physical problem that involves:
    \begin{itemize}
      \item Some number quantities and measurements for them \pause
      \item Some number of implicit relations between them \pause
      \item Some information about the rate of changes of some of the quantities \pause
    \end{itemize}
    and seeks information about the rate of changes of the other variables \pause

    solving the problem via the method of related rates involves: \pause
    \begin{enumerate}
      \item Differentiating the relations between the quantities with respect to the variable of interest \pause
      \item Substituting in the known quantities \pause
      \item Solving the resulting equation(s) for the rate of change in question.
    \end{enumerate}
  \end{block}
\end{frame}

\begin{frame}
  \frametitle{Related Rates Example I}
  For a classic related rates problem:
  Suppose you have a cylindrical tank of water whose height is 10 meters and whose radius is 2 meters,
  with a hole at the bottom from which water drains out.
  At the moment the height of the water is halfway to the bottom, a hole at the bottom of the tank opens and you measure the 
  speed at which the water escapes to be about $\dfrac{1}{2} \dfrac{m^3}{sec}$.
  At what rate will the height of the water decrease? \pause

  Identify information:
  \begin{itemize}
    \item Quantities: \pause volume, height, radius. \pause Radius is constant at 2 meters, volume changes w.r.t. height. \pause Height is at 5 meters.
    \item Relations: \pause $V = \pi r^2 h$. \pause
    \item RoC information: \pause $\dfrac{dV}{dt} = \pause - \dfrac{1}{2} \dfrac{m^3}{sec}$.
    \item RoC in question: \pause $\dfrac{dh}{dt}$.
  \end{itemize}
\end{frame}

\begin{frame}
  \frametitle{Related Rates Example I cont.}
  We then want to differentiate the relation (volume in terms of radius and height) with respect to \pause time to relate the rates of change: \pause

  \begin{gather}
    V = \pi r^2 h \\ \pause
  \dfrac{d}{dt} V = \dfrac{d}{dt} (\pi r^2 h) \\ \pause
  \dfrac{dV}{dt} \dfrac{d}{dV} V = \pi \dfrac{d}{dt} (r^2 h) \\ \pause
  \dfrac{dV}{dt} \cdot 1 = \pi \cdot (\dfrac{d}{dt}(r^2) \cdot h + r^2 \cdot \dfrac{d}{dt}(h)) \\ \pause
  \dfrac{dV}{dt} = \pi \cdot (\dfrac{dr}{dt} \cdot (\dfrac{d}{dr} r^2) \cdot h + r^2 \cdot \dfrac{dh}{dt} \dfrac{d}{dh} h) \\ \pause
  \dfrac{dV}{dt} = \pi \cdot (0 \cdot \_ + r^2 \dfrac{dh}{dt}) \pause \quad \implies \quad \dfrac{dV}{dt} = \pi r^2 \dfrac{dh}{dt} \\ \pause
  -\dfrac{1}{2} = \pi \cdot 2^2 \dfrac{dh}{dt} \pause \quad \implies \quad \dfrac{dh}{dt} = -\dfrac{1}{8\pi} \dfrac{m}{sec}
  \end{gather} \pause

  So the height of the water is decreasing at a rate of $\frac{1}{8\pi}$ meters per second.
\end{frame}

\begin{frame}
  \frametitle{Related Rates Example II}
  Now suppose that the tank is instead an \textit{inverted cone} with the same height and radius. What changes? \pause

  \begin{itemize}
    \item The volume formula changes: \pause $V = \dfrac{1}{3} \pi r^2 h$ \pause
    \item AND \pause the radius also changes over time:
  \end{itemize}

  \vspace{\stretch{1}}

  Giving us the \textit{relation}:
  \begin{gather}
    V = \dfrac{1}{3} \pi r^2 h \\ \pause
    V = \dfrac{1}{3} \pi \left(\dfrac{2}{10} h \right)^2 \cdot h \pause \\
    V = \dfrac{4\pi h^3}{300} \pause = \dfrac{\pi h^3}{75}
  \end{gather}
\end{frame}

\begin{frame}
  \frametitle{Related Rates Example II cont.}
  We would then differentiate this relation w.r.t. $t$ to get:
  \begin{gather}
    \dfrac{d}{dt} V = \dfrac{d}{dt} \dfrac{\pi h^3}{75} \\ \pause
    \dfrac{dV}{dt}  = \dfrac{\pi h^2}{25} \dfrac{dh}{dt} \\ \pause
    \dfrac{dh}{dt} = \dfrac{25}{\pi h^2} \dfrac{dV}{dt} \pause = \dfrac{25}{\pi \cdot 5^2} \cdot \dfrac{-1}{2} = \dfrac{-1}{2\pi} \dfrac{m}{s}
  \end{gather}

  So now the height is decreasing at a rate of $\dfrac{1}{2\pi}$ meters per second.
\end{frame}

\begin{frame}
  \frametitle{Related Rates Example III}
  For a non-tank draining problem, let's instead suppose there
  is a 10 meter long ladder propped up at an angle on a wall that makes a right angle with the ground.
  It's base is currently 6 meters away from the wall, and that base is sliding away from the wall at a rate of two meters per second.
  At that time, how high up along the wall is the ladder, and what is the rate of change of that height? \pause

  \begin{itemize}
    \item Quantities: It is very helpful here to draw a picture! \pause 

      Our quantities are then the distance of the base of the ladder away from the wall $x(t)$ and how high along the wall the ladder is $y(t)$. \pause
    \item Relations: \pause The length of the ladder is constant $\implies$ \pause $x(t)^2 + y(t)^2 = 10^2$. \pause
    \item RoC information: \pause $\dfrac{dx}{dt} = \pause 2 \dfrac{m}{s}$. \pause
    \item Information in question: value of $y(t)$ and $\dfrac{dy}{dt}$.
  \end{itemize}
\end{frame}

\begin{frame}
  \frametitle{Related Rates Example III cont.}
  We differentiate the relation w.r.t time:

  \begin{gather}
    x^2 + y^2 = 10^2 \\ \pause
    \dfrac{d}{dt} (x^2 + y^2) = \dfrac{d}{dt} 10^2 \\ \pause
    \dfrac{d}{dt} (x^2) + \dfrac{d}{dt} (y^2) = 0 \\ \pause
    \dfrac{dx}{dt} \dfrac{d}{dx} x^2 + \dfrac{dy}{dt} \dfrac{d}{dy} y^2 = 0 \\ \pause
    2x \dfrac{dx}{dt} + 2y \dfrac{dy}{dt} = 0 \\ \pause
    2 \cdot 6 \cdot 2 + 2 \cdot y \cdot \dfrac{dy}{dt} = 0 \\ \pause
    \dfrac{dy}{dt} = -\dfrac{12}{y} \dfrac{m}{sec}
  \end{gather}
\end{frame}

\begin{frame}
  \frametitle{Related Rates Example III cont. 2}
  \begin{equation}
    \dfrac{dy}{dt} = - \dfrac{12}{y} \dfrac{m}{sec}
  \end{equation}

  So we still need to determine what $y(t)$ is when $x(t) = 6$. How? \pause

  Use the relation! \pause

  \begin{gather}
    x(t)^2 + y(t)^2 = 10^2 \\  \pause
    6^2 + y(t)^2 = 100 \\ \pause
    y(t)^2 = 100 - 36 \\ \pause
    y(t)^2 = 64 \pause \quad \implies \quad y(t) = \pm 8
  \end{gather}
  and distances must be positive so $y(t) = 8$ when $x(t) = 6$. \pause

  So:
  \begin{equation}
    \dfrac{dy}{dt} = - \dfrac{12}{y} = -\dfrac{12}{8} = -\dfrac{3}{2} \dfrac{m}{sec}.
  \end{equation} \pause

  The height of the ladder against the wall is therefore decreasing at a rate of $\frac{3}{2}$ meters per second.
\end{frame}

\begin{frame}
  \frametitle{More to Come!}

  We'll also be spending next Tuesday going over more examples of how to use 
  the method of related rates to solve different types of problems. \pause

  Generally however, all of these problems will rely on some kind of geometric input, 
  e.g.: \pause

  \begin{multicols}{2}
    \begin{itemize}
      \item Area of square, rectangle, circle, ellipse, triangle
      \item Perimeter of the above (except ellipse)
      \item Angles (e.g. SOHCAHTOA trig) 
      \item Volumes of cubes, rectangular prisms, pyramids, cones, spheres
      \item Surface areas of the cubes, rectangular prisms, spheres
      \item Pythagorean theorem
      \item Similar triangles
    \end{itemize}
  \end{multicols}

  and thus having some way to commit these formulae to memory will be useful.
\end{frame}

\begin{frame}
  \frametitle{Worksheet Time}

  Let's get into groups and work on the worksheet!
\end{frame}

\begin{frame}
  \frametitle{End of Class}

  That's all folks!
\end{frame}
\end{document}
